\documentclass[../../../../uem_mat_mei_q.tex]{subfiles}

\begin{document}
    半径1の円$C_1$に内接する正方形の1つを$S_2$，正方形$S_2$に内接する円を$C_3$，%
    円$C_3$に内接する正三角形の1つを$T_4$，正三角形$T_4$に内接する円を$C_5$とする．%
    以下同様に，自然数$n$について，円$C_{4n+1}$に内接する正方形の1つを$S_{4n+2}$，%
    正方形$S_{4n+2}$に内接する円を$C_{4n+3}$，円$C_{4n+3}$に内接する正三角形の1つを$T_{4n+4}$，%
    正三角形$T_{4n+4}$に内接する円を$C_{4n+5}$とする．

    \begin{myenub}
        \item 円$C_3$の半径は

            \myspaceb%
            $\displaystyle%
                \myfraca{\sqrt{\,\myboxe{ア}\,}}{\myboxe{イ}}%
            $\\
            である．
        \item 正方形$S_6$の面積は

            \myspaceb%
            $\displaystyle%
                \myfracc{\myboxe{ウ}}{\myboxe{エ}}%
            $\\
            である．
        \item 正三角形$T_{16}$の面積は
        
            \myspaceb%
            $\displaystyle%
                3\sqrt{3}\times\left(\myfraca{1}{2}\right)%
                ^{\myboxes{オカ}}%
            $\\
            である．
    \end{myenub}
\end{document}
